\documentclass[10pt,a4paper]{amsart}
\usepackage{amsmath,amssymb,amsthm,graphicx}
\usepackage{verbatim}
\newenvironment{reference}{\comment}{\endcomment}
%\newenvironment{reference}{}{}
\newenvironment{slogan}{\comment}{\endcomment}
\newenvironment{history}{\comment}{\endcomment}
% For commutative diagrams we use Xy-pic
\usepackage[all]{xy}

% We use 2cell for 2-commutative diagrams.
\xyoption{2cell}
% Theorem environments.
%
\theoremstyle{plain}
\newtheorem{theorem}[subsection]{Theorem}
\newtheorem{proposition}[subsection]{Proposition}
\newtheorem{lemma}[subsection]{Lemma}

\theoremstyle{definition}
\newtheorem{definition}[subsection]{Definition}
\newtheorem{example}[subsection]{Example}
\newtheorem{exercise}[subsection]{Exercise}
\newtheorem{situation}[subsection]{Situation}

\theoremstyle{remark}
\newtheorem{remark}[subsection]{Remark}
\newtheorem{remarks}[subsection]{Remarks}

\numberwithin{equation}{subsection}

% Macros
%
\def\lim{\mathop{\mathrm{lim}}\nolimits}
\def\colim{\mathop{\mathrm{colim}}\nolimits}
\def\Spec{\mathop{\mathrm{Spec}}}
\def\Hom{\mathop{\mathrm{Hom}}\nolimits}
\def\Ext{\mathop{\mathrm{Ext}}\nolimits}
\def\SheafHom{\mathop{\mathcal{H}\!\mathit{om}}\nolimits}
\def\SheafExt{\mathop{\mathcal{E}\!\mathit{xt}}\nolimits}
\def\Sch{\mathit{Sch}}
\def\Mor{\mathop{\mathrm{Mor}}\nolimits}
\def\Ob{\mathop{\mathrm{Ob}}\nolimits}
\def\Sh{\mathop{\mathit{Sh}}\nolimits}
\def\NL{\mathop{N\!L}\nolimits}
\def\CH{\mathop{\mathrm{CH}}\nolimits}
\def\proetale{{pro\text{-}\acute{e}tale}}
\def\etale{{\acute{e}tale}}
\def\QCoh{\mathit{QCoh}}
\def\Ker{\mathop{\mathrm{Ker}}}
\def\Im{\mathop{\mathrm{Im}}}
\def\Coker{\mathop{\mathrm{Coker}}}
\def\Coim{\mathop{\mathrm{Coim}}}

% Boxtimes
%
\DeclareMathSymbol{\boxtimes}{\mathbin}{AMSa}{"02}

%
% Macros for moduli stacks/spaces
%
\def\QCohstack{\mathcal{QC}\!\mathit{oh}}
\def\Cohstack{\mathcal{C}\!\mathit{oh}}
\def\Spacesstack{\mathcal{S}\!\mathit{paces}}
\def\Quotfunctor{\mathrm{Quot}}
\def\Hilbfunctor{\mathrm{Hilb}}
\def\Curvesstack{\mathcal{C}\!\mathit{urves}}
\def\Polarizedstack{\mathcal{P}\!\mathit{olarized}}
\def\Complexesstack{\mathcal{C}\!\mathit{omplexes}}
% \Pic is the operator that assigns to X its picard group, usage \Pic(X)
% \Picardstack_{X/B} denotes the Picard stack of X over B
% \Picardfunctor_{X/B} denotes the Picard functor of X over B
\def\Pic{\mathop{\mathrm{Pic}}\nolimits}
\def\Picardstack{\mathcal{P}\!\mathit{ic}}
\def\Picardfunctor{\mathrm{Pic}}
\def\Deformationcategory{\mathcal{D}\!\mathit{ef}}
\usepackage{cleveref}
\crefname{originallemma}{original lemma}{original lemmas}
\Crefname{originallemma}{Original lemma}{Original lemmas}
\makeatletter
\newlabel{lemma-power-ideal-kills-sheaf}{{07UK}{source}{}{sourcepointer-1}{}}
\newlabel{lemma-power-ideal-kills-sheaf@cref}{{[originallemma][0][]07UK}{[1][0][]source}{}{}{}}
\makeatother
\crefname{originallemma}{original lemma}{original lemmas}
\Crefname{originallemma}{Original lemma}{Original lemmas}
\makeatletter
\newlabel{lemma-prepare-filter-irreducible}{{07UQ}{source}{}{sourcepointer-2}{}}
\newlabel{lemma-prepare-filter-irreducible@cref}{{[originallemma][0][]07UQ}{[1][0][]source}{}{}{}}
\makeatother
\crefname{originallemma}{original lemma}{original lemmas}
\Crefname{originallemma}{Original lemma}{Original lemmas}
\makeatletter
\newlabel{lemma-prepare-filter-support}{{07UP}{source}{}{sourcepointer-3}{}}
\newlabel{lemma-prepare-filter-support@cref}{{[originallemma][0][]07UP}{[1][0][]source}{}{}{}}
\makeatother
\crefname{originallemma}{original lemma}{original lemmas}
\Crefname{originallemma}{Original lemma}{Original lemmas}
\makeatletter
\newlabel{spaces-morphisms-lemma-i-star-equivalence}{{04CJ}{source}{}{sourcepointer-4}{}}
\newlabel{spaces-morphisms-lemma-i-star-equivalence@cref}{{[originallemma][0][]04CJ}{[1][0][]source}{}{}{}}
\makeatother
\crefname{originallemma}{original lemma}{original lemmas}
\Crefname{originallemma}{Original lemma}{Original lemmas}
\makeatletter
\newlabel{spaces-properties-lemma-Noetherian-topology}{{04ZF}{source}{}{sourcepointer-5}{}}
\newlabel{spaces-properties-lemma-Noetherian-topology@cref}{{[originallemma][0][]04ZF}{[1][0][]source}{}{}{}}
\makeatother
\crefname{originallemma}{original lemma}{original lemmas}
\Crefname{originallemma}{Original lemma}{Original lemmas}
\makeatletter
\newlabel{spaces-properties-lemma-reduced-closed-subspace}{{03IQ}{source}{}{sourcepointer-6}{}}
\newlabel{spaces-properties-lemma-reduced-closed-subspace@cref}{{[originallemma][0][]03IQ}{[1][0][]source}{}{}{}}
\makeatother
\title{Every coherent sheaf on a Noetherian algebraic space is filtered with subquotients pushed forward from ideal sheaves on irreducible supports}
\author{The Stacks Project Authors}
\date{stacks-project:spaces-cohomology.tex:local-snapshot-20261012; GNU FDL 1.2 or later}
\begin{document}
\maketitle
\noindent\textit{Editorial note.} All delivered mathematical text is the author's own, copied from the pinned Stacks Project chapter file spaces-cohomology.tex; no statement or proof step was rewritten, reordered or removed.\par\medskip
\noindent\textit{Editorial note.} Preamble provenance: the pinned chapter file opens with \textbackslash{}input\{preamble\} and carries no preamble text of its own, so the project's single Stacks-wide preamble staged at sources/stacks-spaces-cohomology/preamble.tex is excerpted verbatim as three spans, lines 8-12 (verbatim package and the project's own reference/slogan/history display environments), lines 14-18 (Xy-pic setup) and lines 154-217 (theorem declarations and mathematical macros). The wrapper class and the amsmath/amssymb/amsthm/graphicx base packages are supplied by the assembler. No macro was invented or redefined.\par\medskip
\noindent\textit{Editorial note.} Same-chapter pointers are located at their \textbackslash{}label lines inside the pinned chapter file; cross-chapter pointers are located in the project's own tag list staged in the snapshot (sources/stacks-chow/tags.txt), at the verbatim lines carrying their canonical tags.\par\medskip
\noindent\textit{Allowed standard prerequisites (editorial).}\begin{itemize}
\item Splitting a coherent sheaf along a union of two strictly smaller closed supports. Same chapter spaces-cohomology.tex, label lemma-prepare-filter-support (invoked at proof line 2565). It produces the short exact sequence \$0 \textbackslash{}to \textbackslash{}mathcal\{G\}\_1 \textbackslash{}to \textbackslash{}mathcal\{F\} \textbackslash{}to \textbackslash{}mathcal\{G\}\_2 \textbackslash{}to 0\$ with supports in the two smaller closed pieces, which turns the reducible case into a contradiction with minimality.
\item A power of the ideal of the support annihilates the sheaf. Same chapter spaces-cohomology.tex, label lemma-power-ideal-kills-sheaf (invoked at proof line 2584). It gives \$\textbackslash{}mathcal\{J\}\textasciicircum{}n\textbackslash{}mathcal\{F\} = 0\$, so that the filtration by the powers of \$\textbackslash{}mathcal\{J\}\$ reduces the assertion to the case where \$\textbackslash{}mathcal\{J\}\$ annihilates \$\textbackslash{}mathcal\{F\}\$.
\item The scheme theoretic support, so that the sheaf is a pushforward from the support. Morphisms of Spaces, label spaces-morphisms-lemma-i-star-equivalence (invoked at proof line 2600). It identifies the scheme theoretic support of \$\textbackslash{}mathcal\{F\}\$ with \$T\$, so that \$\textbackslash{}mathcal\{F\}\$ is the pushforward of a coherent module on \$T\$ and the preparation lemma for the irreducible case applies.
\item Preparation of a coherent module on a reduced irreducible space by a pushed forward ideal power. Same chapter spaces-cohomology.tex, label lemma-prepare-filter-irreducible (invoked at proof line 2601). It gives the short exact sequence \$0 \textbackslash{}to i\_*(\textbackslash{}mathcal\{I\}\textasciicircum{}\{\textbackslash{}oplus r\}) \textbackslash{}to \textbackslash{}mathcal\{F\} \textbackslash{}to \textbackslash{}mathcal\{Q\} \textbackslash{}to 0\$ with \$\textbackslash{}mathcal\{Q\}\$ supported in a proper closed subset of \$T\$, whose subquotients have the required shape.
\item Noetherianity of the underlying topological space, and the reduced induced structure on a closed subset. Properties of Spaces, labels spaces-properties-lemma-Noetherian-topology (invoked at proof line 2554) and spaces-properties-lemma-reduced-closed-subspace (invoked at proof line 2582). The first supplies the minimal element of the collection of counterexample supports; the second provides the ideal sheaf defining the reduced structure on that support.
\end{itemize}\medskip
\ifdefined\hypertarget\hypertarget{sourcepointer-1}{}\fi
\noindent\textit{Source reference.} Original lemma 07UK: sources/stacks-spaces-cohomology/pinned-original-source.tex, source lines 2339--2339. Navigation pointer for \textbackslash{}ref\{lemma-power-ideal-kills-sheaf\} used inside the counted span (line 2584); the argument that consumes it is retained in the span.\par\medskip
\ifdefined\hypertarget\hypertarget{sourcepointer-2}{}\fi
\noindent\textit{Source reference.} Original lemma 07UQ: sources/stacks-spaces-cohomology/pinned-original-source.tex, source lines 2478--2478. Navigation pointer for \textbackslash{}ref\{lemma-prepare-filter-irreducible\} used inside the counted span (line 2601); the argument that consumes it is retained in the span.\par\medskip
\ifdefined\hypertarget\hypertarget{sourcepointer-3}{}\fi
\noindent\textit{Source reference.} Original lemma 07UP: sources/stacks-spaces-cohomology/pinned-original-source.tex, source lines 2435--2435. Navigation pointer for \textbackslash{}ref\{lemma-prepare-filter-support\} used inside the counted span (line 2565); the argument that consumes it is retained in the span.\par\medskip
\ifdefined\hypertarget\hypertarget{sourcepointer-4}{}\fi
\noindent\textit{Source reference.} Original lemma 04CJ: sources/stacks-chow/tags.txt, source lines 5693--5693. Navigation pointer for \textbackslash{}ref\{spaces-morphisms-lemma-i-star-equivalence\} used inside the counted span (line 2600); the argument that consumes it is retained in the span.\par\medskip
\ifdefined\hypertarget\hypertarget{sourcepointer-5}{}\fi
\noindent\textit{Source reference.} Original lemma 04ZF: sources/stacks-chow/tags.txt, source lines 6463--6463. Navigation pointer for \textbackslash{}ref\{spaces-properties-lemma-Noetherian-topology\} used inside the counted span (line 2555); the argument that consumes it is retained in the span.\par\medskip
\ifdefined\hypertarget\hypertarget{sourcepointer-6}{}\fi
\noindent\textit{Source reference.} Original lemma 03IQ: sources/stacks-chow/tags.txt, source lines 4612--4612. Navigation pointer for \textbackslash{}ref\{spaces-properties-lemma-reduced-closed-subspace\} used inside the counted span (line 2583); the argument that consumes it is retained in the span.\par\medskip
\begin{lemma}
\label{lemma-coherent-filter}
Let $S$ be a scheme. Let $X$ be a Noetherian algebraic space over $S$.
Let $\mathcal{F}$ be a coherent sheaf on $X$. There exists a filtration
$$
0 = \mathcal{F}_0 \subset \mathcal{F}_1 \subset
\ldots \subset \mathcal{F}_m = \mathcal{F}
$$
by coherent subsheaves such that for each $j = 1, \ldots, m$
there exists a reduced closed subspace $Z_j \subset X$ with $|Z_j|$
irreducible and a sheaf of ideals $\mathcal{I}_j \subset \mathcal{O}_{Z_j}$
such that
$$
\mathcal{F}_j/\mathcal{F}_{j - 1}
\cong (Z_j \to X)_* \mathcal{I}_j
$$
\end{lemma}

\begin{proof}
Consider the collection
$$
\mathcal{T} =
\left\{
\begin{matrix}
T \subset |X|
\text{ closed such that there exists a coherent sheaf }
\mathcal{F} \\
\text{ with }
\text{Supp}(\mathcal{F}) = T
\text{ for which the lemma is wrong}
\end{matrix}
\right\}
$$
We are trying to show that $\mathcal{T}$ is empty. If not, then
because $|X|$ is Noetherian (Properties of Spaces, Lemma
\ref{spaces-properties-lemma-Noetherian-topology})
we can choose a minimal element $T \in \mathcal{T}$. This means that
there exists a coherent sheaf $\mathcal{F}$ on $X$ whose support is $T$
and for which the lemma does not hold. Clearly $T \not = \emptyset$ since
the only sheaf whose support is empty is the zero sheaf for which the
lemma does hold (with $m = 0$).

\medskip\noindent
If $T$ is not irreducible, then we can write $T = Z_1 \cup Z_2$
with $Z_1, Z_2$ closed and strictly smaller than $T$.
Then we can apply Lemma \ref{lemma-prepare-filter-support}
to get a short exact sequence of coherent sheaves
$$
0 \to
\mathcal{G}_1 \to
\mathcal{F} \to
\mathcal{G}_2 \to 0
$$
with $\text{Supp}(\mathcal{G}_i) \subset Z_i$. By minimality of
$T$ each of $\mathcal{G}_i$ has a filtration as in the statement
of the lemma. By considering the induced filtration on $\mathcal{F}$
we arrive at a contradiction. Hence we conclude
that $T$ is irreducible.

\medskip\noindent
Suppose $T$ is irreducible. Let $\mathcal{J}$ be the sheaf of ideals
defining the reduced induced closed subspace structure on $T$,
see Properties of Spaces, Lemma
\ref{spaces-properties-lemma-reduced-closed-subspace}.
By Lemma \ref{lemma-power-ideal-kills-sheaf} we see there exists
an $n \geq 0$ such that $\mathcal{J}^n\mathcal{F} = 0$. Hence we obtain
a filtration
$$
0 = \mathcal{I}^n\mathcal{F} \subset \mathcal{I}^{n - 1}\mathcal{F}
\subset \ldots \subset \mathcal{I}\mathcal{F} \subset \mathcal{F}
$$
each of whose successive subquotients is annihilated by $\mathcal{J}$.
Hence if each of these subquotients has a filtration as in the statement
of the lemma then also $\mathcal{F}$ does. In other words we may
assume that $\mathcal{J}$ does annihilate $\mathcal{F}$.

\medskip\noindent
Assume $T$ is irreducible and $\mathcal{J}\mathcal{F} = 0$ where
$\mathcal{J}$ is as above. Then the scheme theoretic support of
$\mathcal{F}$ is $T$, see
Morphisms of Spaces, Lemma \ref{spaces-morphisms-lemma-i-star-equivalence}.
Hence we can apply Lemma \ref{lemma-prepare-filter-irreducible}.
This gives a short exact sequence
$$
0 \to
i_*(\mathcal{I}^{\oplus r}) \to
\mathcal{F} \to
\mathcal{Q} \to 0
$$
where the support of $\mathcal{Q}$ is a proper closed subset of $T$.
Hence we see that $\mathcal{Q}$ has a filtration of the desired type
by minimality of $T$. But then clearly $\mathcal{F}$ does too, which is
our final contradiction.
\end{proof}

\end{document}
